\documentclass[../../../include/open-logic-section]{subfiles}

\begin{document}

\olfileid{sth}{cardinals}{hp}
\olsection{ملحق: مبدأ هيوم}

مر وصف مبدأ كانتور في~\olref[cp]{sec}، وصورته:
\begin{align*}
	\card{A} = \card{B} & \text{ إذا وفقط إذا كان } A \approx B.
\intertext{وهذا شبيه جداً بما يسمى اليوم
\emph{مبدأ هيوم}، الذي يقول:}
	\fregenum{x} {F(x)} = \fregenum{x}{G(x)} & \text{ إذا وفقط إذا كان } F \sim G
\end{align*}
ويراد بـ«$F \sim G$» أن عدد الـ$F$ات مساو تمامًا لعدد الـ$G$ات؛ أي
إنه يمكن إقامة تقابل بين الـ$F$ات والـ$G$ات، وصورته:
\begin{align*}
	\exists R(&\forall v\forall y(Rvy \lif (Fv \land Gy)) \land {}\\
		&\forall v(Fv \lif \lexists![y][Rvy]) \land {}\\
		&\forall y(Gy \lif \lexists![v][Rvy]))
\end{align*}
غير أن مبدأ هيوم يخالف مبدأ كانتور في النوع. ففي كانتور المتغيران
«$A$» و«$B$» حدان من الرتبة الأولى، يدلان على \emph{مجموعتين}. وفي
هيوم لا تكون «$F$» و«$G$» و«$R$» حدودًا من الرتبة الأولى، بل ترد في
\emph{موضع المحمول}، ولعلها تدل على \emph{خصائص}. ولذلك يكون شرحه
بالعربية: عدد الـ$F$ات يساوي عدد الـ$G$ات إذا وفقط إذا أمكن إقامة
تقابل بين الـ$F$ات والـ$G$ات. وإنما نسب إلى هيوم لقوله:
\begin{quote}
  حين يُقرن عددان بحيث تقابل دائماً كل وحدة من أحدهما وحدة من الآخر،
  نحكم بأنهما متساويان.
  \citep[Pt.III Bk.1 \S1]{Hume1740}
\end{quote}
ثم جعل~\citet[\S63]{Frege1884} لهذا المبدأ مكانه البارز في المنطق
الرياضي الحديث؛ إذ ساق كلام هيوم قبل أن يخط، في الحقيقة، معالم ما
سميناه مبدأ هيوم.

وقارن البنية هنا بالقانون الأساسي الخامس، الذي صيغ في
\olref[story][blv]{sec} هكذا:
\[
	\fregeext{x}{F(x)} = \fregeext{x}{G(x)} \text{إذا وفقط إذا كان } \lforall[x][(F(x) \liff G(x))].
\]
وفي هذه المقارنة ما يستحق البيان.

فكل مبدأ من قبيل مبدأ هيوم أو القانون الأساسي الخامس يحتمل قراءتين:
قراءة \emph{حملية} وأخرى \emph{غير حملية}؛ راجع
\olref[story][predicative]{sec}. فعلى القراءة غير الحملية للقانون
الأساسي الخامس، يكون الكائن~$\fregeext{x}{F(x)}$، لكل~$F$، داخل مجال
التكميم الذي صيغ به القانون نفسه. وعلى مثلها في مبدأ هيوم، يكون
$\fregenum{x}{F(x)}$، لكل~$F$، داخل مجال التكميم الذي صيغ به المبدأ.
وأما على القراءة \emph{الحملية}، فالكائنان
$\fregeext{x}{F(x)}$ و$\fregenum{x}{F(x)}$ من مجال \emph{آخر}.

والقانون الأساسي الخامس، إذا قرئ قراءة غير حملية، يفضي مع الاستيعاب
الساذج إلى عدم الاتساق؛ وانظر~\olref[story][blv]{sec}. وكما نفعل في
الاستيعاب الساذج، يمكن جعله متسقًا بقراءته \emph{حمليًا}، لكنه على
الأرجح لا يؤدي كل الغرض الذي أردناه منه.

أما مبدأ هيوم فيصح، \emph{مع الاتساق}، أن يقرأ قراءة غير حملية، ويكون
عندئذ بالغ القوة.

ولبيان ذلك، خذ المحمول «$x \neq x$»، وهو لا يصدق في شيء. يمنحنا مبدأ
هيوم الكائن~$\# x( x\neq x)$، ولنا أن نعامله بوصفه العدد~$0$. وعلى
الفهم \emph{غير الحملي}، ولا يتم هذا إلا عليه، يقع هذا الكائن~$0$
في مجال التكميم الأصلي. فيصح تطبيق مبدأ هيوم على المحمول «$x = 0$»،
فنحصل على~$\#x (x = 0)$، ونعامله بوصفه العدد~$1$. ويقتضي المبدأ كذلك
$0 \neq 1$؛ إذ لا تقابل بين الكائنات غير المطابقة لأنفسها والكائنات
المطابقة لـ$0$: الأولى لا شيء، والثانية واحد. وإذا واصلنا العمل على
الوجه غير الحملي، كان~$1$ داخل مجال التكميم الأصلي أيضًا. فنطبق المبدأ
على المحمول «$(x = 0 \lor x = 1)$»، فنحصل على الكائن
$\#x(x = 0 \lor x = 1)$، ونعامله بوصفه العدد~$2$. ثم نبين
$0\neq 2$ و$1 \neq 2$، ونمضي على هذا القياس.

فمبدأ هيوم، مأخوذًا على الوجه غير الحملي، يستلزم وجود \emph{عدد غير
متناه من الكائنات}. وهذا ما شجع \emph{المناطقة الفريغيين الجدد} على
جعله أساسًا للحساب.

لكن فرغه \emph{نفسه} لم يؤسس الحساب عليه، بل أثبته من تعريف صريح:
عرّف~$\fregenum{x}{F(x)}$ بأنه امتداد المفهوم~$F \sim \Phi$.
وباصطلاح حديث قد تكتب المحاولة هكذا:
$\fregenum{x}{F(x)} = \Setabs{G}{F \sim G}$؛ غير أنها تردنا إلى
إشكالات الاستيعاب الساذج.

\end{document}
